\documentclass[11pt,a4paper]{article}
\usepackage[margin=2.5cm]{geometry}
\usepackage{amsmath,amssymb}
\usepackage{booktabs}
\usepackage{hyperref}

\newcommand{\bM}{\boldsymbol{M}}
\newcommand{\bp}{\boldsymbol{p}}

\begin{document}

\title{Calculation of RR Photon Spectrum including Radiative Cascades in \texttt{hydrocal/src/RRphotons.cxx}}
\author{AI generated and edited by S. Schippers}
\date{\today}
\maketitle

\section{Problem}

When a free electron is captured by a hydrogenic ion via radiative
recombination (RR), two kinds of photons are emitted:
\begin{enumerate}
\item \textbf{Free-bound (RR) photon}: emitted during the capture of
  a free electron with kinetic energy $E_{\rm kin}$ and angular
  momentum $l'$ into a bound state $|n,l\rangle$ with
  $l = l' \pm 1$.
\item \textbf{Cascade photons}: emitted as the excited ion decays
  through a chain of dipole transitions down to the ground state.
\end{enumerate}
The goal is to compute a complete line list---photon energies,
intensities, and natural linewidths---that can later be convolved with
a Gaussian instrumental profile using the existing \texttt{sum\_of\_peaks}
tool.

\section{Free-Bound (RR) Photons}

\subsection{Photon Energy}

The RR photon energy in the ion rest frame is determined by energy
conservation.  For high-$z$ ions, relativistic (Dirac) energies are
used instead of the non-relativistic Bohr formula.  The Dirac binding
energy of a hydrogenic state $|n,l\rangle$ is computed by averaging over
the two fine-structure levels $j = l - 1/2$ and $j = l + 1/2$ with
statistical weights $(2j{+}1)$:
\begin{equation}
  E_{\rm Dirac}(n,l) = \frac{2l\; E(n,l,j{=}l{-}\tfrac{1}{2})
    + 2(l{+}1)\; E(n,l,j{=}l{+}\tfrac{1}{2})}{4l + 2}
  \label{eq:Dirac}
\end{equation}
(for $l = 0$ only $j = 1/2$ exists).  Each component is
given by:
\begin{equation}
  E(n,l,j) = m_e c^2
  \left(\frac{1}{\sqrt{1 + \kappa^2}} - 1\right)
  \label{eq:Diraccomp}
\end{equation}
where
\begin{equation}
  \kappa = \frac{z\,\alpha}{n - k + \sqrt{k^2 - (z\,\alpha)^2}}
  \label{eq:kappa}
\end{equation}
with $\alpha$ the fine-structure constant, $k = j + 1/2$, and the
result averaged over the two allowed $j$ values for $l > 0$.  The RR
photon energy is then:
\begin{equation}
  E_{\rm RR}(n,l) = E_{\rm kin} - E_{\rm Dirac}(n,l)
  \label{eq:ERR}
\end{equation}
where $E_{\rm Dirac}(n,l) < 0$ is the (negative) binding energy.  The
energy now depends on both $n$ and~$l$ due to fine-structure splitting,
lifting the $l$-degeneracy of the non-relativistic formula.

\subsection{Line Strength}

The RR cross section $\sigma_{\rm RR}(n,l,E_{\rm kin})$ is obtained
from the quantum-mechanical calculation implemented in
\texttt{sigmarrqm()}~\cite{Schippers2002c}.  The line strength of each
RR photon equals the cross section for capture into $(n,l)$:
\begin{equation}
  S_{\rm RR}(n,l) = \sigma_{\rm RR}(n,l,E_{\rm kin})
  \qquad [\text{cm}^2]
  \label{eq:SRR}
\end{equation}
The output lists one entry per $(n,l)$ state (not grouped by $n$), so
the $l$-distribution of the initial capture is preserved.

\subsection{Natural Linewidth}

The RR photon has no natural broadening from the final state lifetime
(the photon energy is fixed by energy conservation), so the natural
width is set to zero.  Any broadening comes from the experimental
electron energy spread and is applied later via convolution.

\subsection{Label Convention}

Each RR line is labelled $(E,l{-}1)/(E,l{+}1) \to (n,l)$ for $l > 0$,
or $(E,1) \to (n,0)$ for $l = 0$, indicating the allowed free-electron
angular momenta consistent with the dipole selection rule $\Delta l =
\pm 1$.  The letter ``$E$'' represents the continuum electron.  For a
relativistic target state ($n \le n_{\rm rel}$) the label additionally
carries $2j$: $(E,l{-}1)/(E,l{+}1) \to (n,l,2j)$; the continuum
partial waves then correspond to the E1-coupled channels
$\kappa_c = \kappa \pm 1$ and $\kappa_c = -\kappa$.

\section{Cascade Photons}

\subsection{Rate Matrix}

The radiative cascade is governed by the coupled rate equations
\begin{equation}
  \frac{d p_i}{dt} = \sum_j M_{ij}\, p_j
  \label{eq:rate}
\end{equation}
where $p_i$ is the population of state $i$ (index $i = n(n{-}1)/2 + l$)
and $\bM$ is the rate matrix:
\begin{equation}
  M_{ii} = -\frac{1}{\tau_i}\,, \qquad
  M_{ji} = \frac{B(i \to j)}{\tau_i} \quad (j \ne i)
  \label{eq:M}
\end{equation}
with $\tau_i$ the radiative lifetime of state~$i$ and $B(i \to j)$ the
branching ratio for the transition $i \to j$ (the sum of the E1, E2,
and M1 rates; parity keeps E1 separate from E2/M1).  The ground state
$(1,0)$ has $\tau = \infty$ (no decay), so its column in $\bM$ is zero.
Among relativistic states ($n \le n_{\rm rel}$) each fine-structure
level $j$ is a separate state, and E1 as well as E2/M1 transitions
enter the rate matrix as individual entries $M_{ji} = A(i\to j)$.

\subsection{Time-Integrated Populations}

The total photon yield from each transition requires the time-integrated
population of the upper state:
\begin{equation}
  \bar{p}_i = \int_0^\infty p_i(t)\, dt
  \label{eq:pb}
\end{equation}
Since the ground state makes $\bM$ singular, we restrict to the
excited-state submatrix $\bM_{\rm exc}$ of dimension
$N_{\rm exc} = n_{\rm max}(n_{\rm max}{+}1)/2 - 1$ (all states except
$(1,0)$).  Then:
\begin{equation}
  \bar{\bp}_{\rm exc} = -\bM_{\rm exc}^{-1}\, \bp_0
  \label{eq:barp}
\end{equation}
where $\bp_0$ is the initial population vector (RR cross sections) and
$\bM_{\rm exc}$ is invertible because all its eigenvalues have negative
real parts.

The linear system
\begin{equation}
  \bM_{\rm exc}\, \bar{\bp}_{\rm exc} = -\bp_0
  \label{eq:lin}
\end{equation}
is solved by Gaussian elimination with partial pivoting
(\texttt{matrixSolve()}).

\subsection{Photon Energy}

For each dipole-allowed transition $(n_1,l_1) \to (n_2,l_2)$ with
$n_2 < n_1$ and $l_2 = l_1 \pm 1$:
\begin{equation}
  E_{\gamma} = E_{\rm Dirac}(n_1,l_1) - E_{\rm Dirac}(n_2,l_2)
  \label{eq:Ecasc}
\end{equation}
where $E_{\rm Dirac}(n,l)$ is given by Eq.~(\ref{eq:Dirac}).  For
high-$z$ ions the fine-structure splitting of both states is taken into
account, so cascade photon energies differ from the non-relativistic
Balmer/Ritz formula.

\subsection{Line Strength}

The total number of photons emitted in transition $i \to j$,
integrated over all time, is:
\begin{equation}
  S(i \to j) = A(i \to j)\; \bar{p}_i
  \label{eq:Scasc}
\end{equation}
where $A(i \to j) = \texttt{trans}(n_1,l_1,n_2,l_2)$ is the
Einstein~$A$ coefficient (transition rate in~s$^{-1}$) and
$\bar{p}_i$ is the time-integrated population of the upper state in
units of s$\cdot$cm$^2$.  The product has units of cm$^2$.

\subsection{Natural Linewidth}

The natural Lorentzian FWHM of a spectral line is determined by the
lifetimes of both the upper and lower states:
\begin{equation}
  \Gamma_{ij} = \hbar \left(\frac{1}{\tau_i} + \frac{1}{\tau_j}\right)
  \label{eq:width}
\end{equation}
For transitions to the ground state ($n_2 = 1$), $\tau_j \to \infty$
and the width reduces to $\Gamma = \hbar / \tau_i$.

\subsection{Label Convention}

Each cascade line is labelled $(n_1,l_1) \to (n_2,l_2)$ with the
actual quantum numbers.  For relativistic states ($n \le n_{\rm rel}$)
the label is $j$-resolved, $(n_1,l_1,2j_1) \to (n_2,l_2,2j_2)$, and
the rate is $A_{\text{total}} = A_{\text{E1}}$ or
$A_{\text{E2}}+A_{\text{M1}}$ depending on the
selection rules.  E2/M1 lines (e.g.\ the fine-structure transition
$2p_{3/2}\to2p_{1/2}$) therefore appear as individual spectral lines
at their characteristic energy.

\section{Relativistic treatment for $n \le n_{\rm rel}$}

For high-$z$ ions (or when $n_{\rm rel} > 0$) the states with
$n \le n_{\rm rel}$ are described by Dirac hydrogenic spinors with a
definite total angular momentum $j$ (encoded in $\kappa$), while
higher-$n$ states keep the non-relativistic treatment of the previous
sections.  For such states:
\begin{itemize}
\item The binding energy is the Dirac value $E(n,\kappa)$; the two
  fine-structure components of an $(n,l)$ shell are treated as
  individual states and labelled $(n,l,2j)$ (e.g.\ $(2,1,1)$ for
  $2p_{1/2}$ and $(2,1,3)$ for $2p_{3/2}$).
\item The RR cross section for capture into $(n,\kappa)$ is obtained
  from the relativistic Milne relation, summed over the E1-coupled
  continuum channels $\kappa_c = \kappa\pm1$ and $\kappa_c = -\kappa$.
  The direct-RR label is $(E,l{-}1)/(E,l{+}1) \to (n,l,2j)$.
\item Cascade transitions between two relativistic states are
  $j$-resolved: the rate is $A_{\text{total}}(n_i,\kappa_i \to
  n_j,\kappa_j)$ (sum of E1, E2, and M1; E1 never coincides with
  E2/M1 because of parity, so the sum equals the allowed multipole),
  the photon energy
  is the difference of the two Dirac binding energies, and the E2/M1
  transitions enter both the rate matrix and the line list as
  individual entries.
\item Decays from a non-relativistic upper state into a relativistic
  fine-structure doublet are distributed by statistical weight
  $(2j{+}1)/[2(2l{+}1)]$.
\end{itemize}

The rate matrix, population solution, and line construction are
otherwise unchanged; every transition in the final line list connects
states with definite $j$ wherever the lower state is relativistic.

\section{Output Format}

The output file (\texttt{*.RRlines}) is a columnar text file compatible
with the existing \texttt{read\_peak()} / \texttt{sum\_of\_peaks()}
infrastructure.  Each data line is preceded by a \texttt{\#\#\#} comment
line carrying the transition label:
\begin{verbatim}
#####################################################################
### Photon spectrum resulting from the radiative cascade following
### radiative recombination (RR lines + cascade lines)
###
###  hydrocal SVN revision : 2089:2090M
###               filename : Pb82RRphotons.RRlines
###  number of spec. lines : 2680
###      start date & time : Fri Jul 31 11:28:21 2026
###    nuclear charge Zeff : 82
###   electron energy (eV) : 1e-05
###                   nmax : 20
### min photon energy (eV) : 0
###  relativistic for n<=  : 0
#####################################################################
#### energy (eV) [1]   cross_section (cm^2) [2]   line-width (eV) [3]
###-----------------------------------------------------------------
### (20,0)->(19,1)
  2.34962616e+01   6.17478755e-18   1.60586110e-01
### (20,1)->(19,2)
  2.36837825e+01   1.11557824e-18   1.78156903e-01
  ...
\end{verbatim}
Lines starting with \texttt{\#} are skipped by \texttt{read\_peak()}
via the \texttt{skip\_comment\_lines()} function.  The three numeric
columns are:
\begin{enumerate}
\item \textbf{Energy} (eV): photon energy in the ion rest frame.
\item \textbf{Strength} (cm$^2$): integrated line intensity
  (cross section for RR lines; transition rate $\times$ time-integrated
  population for cascade lines).
\item \textbf{Width} (eV): natural Lorentzian FWHM (zero for RR lines).
\end{enumerate}
Lines are sorted by photon energy.

\section{Convolution}

The line list can be convolved with a Gaussian instrumental profile using
\texttt{sum\_of\_peaks} (mode~2~= Lorentz).  The tool reads the
\texttt{*.RRlines} file, prompts for a Gaussian FWHM, and produces a
broadened spectrum as a function of photon energy.  Voigt profiles
(Lorentzian $\otimes$ Gaussian) are obtained automatically when the
input contains non-zero Lorentzian widths.

\section{Implementation Notes}

\begin{itemize}
\item All cascade computations (fully non-relativistic and mixed
  relativistic/non-relativistic) are handled by a single code path
  (\texttt{cascade\_spectrum\_rel()} in \texttt{RRphotons.cxx}).  For
  $n_{\rm rel}=0$ all states are treated non-relativistically; for
  $n_{\rm rel}>0$ the states with $n \le n_{\rm rel}$ are treated with the
  Dirac rates and the higher-$n$ states with the non-relativistic rates.
\item The rate matrix $\bM$ carries the individual
  relativistic E1, E2 and M1 transition rates (from
  \texttt{dirac\_total\_transrate()}) and the non-relativistic dipole
  branching ratios (from \texttt{RADRATE::branch()}) as its off-diagonal
  entries.  Decays from non-relativistic upper states require
  $\Delta l = \pm 1$ (dipole selection rule); decays between relativistic
  states are restricted by the Dirac parity and angular-momentum
  selection rules.
\item The \texttt{RADRATE} class provides transition rates, branching
  ratios, and lifetimes computed from hydrogenic oscillator strengths.
  The rates are in~s$^{-1}$ and the lifetimes in~s.
\item The linear system $\bM_{\rm exc}\, \bar{\bp} = -\bp_0$ is solved
  by Gaussian elimination with partial pivoting
  (\texttt{matrixSolve()} in \texttt{hydromath.cxx}).
\item For zero electron energy ($E_{\rm kin} < 10^{-10}$~eV), the RR
  cross section is evaluated at $E = 10^{-12}$~eV to avoid the $1/E$
  singularity of the Wigner threshold law.
\item The \texttt{calc\_spectrum()} function is accessible from the
  hydrocal menu as RR cross sections group, option~4.
\end{itemize}

\begin{thebibliography}{1}
\bibitem{Schippers2002c}
   S. Schippers, Experimente zur Photorekombination Atomarer Ionen an Schwerionenspeicherringen (\href{https://doi.org/10.22029/jlupub-15034}{Habilitation Thesis}), Justus-Liebig-Universit\"at Gie{\ss}en, Germany, 2002, Chapter~3.
\end{thebibliography}

\end{document}
